% problems4.tex — Problem Set and Answers for Chapter 4

\begin{problemset}

\subsection*{Analytical Exercises}

\exercise[1]
(Data matrix representation of 3SLS)\quad
In the 3SLS model, the set of instruments is the same across equations.
Let $K$ be the number of common instruments. Define:
\[
\underset{(n \times K)}{\mathbf{X}}
= \begin{bmatrix} \mathbf{x}_1' \\ \vdots \\ \mathbf{x}_n' \end{bmatrix}
= \begin{bmatrix}
x_{11} & \cdots & x_{1K} \\
\vdots & & \vdots \\
x_{n1} & \cdots & x_{nK}
\end{bmatrix},
\]
\[
\underset{(n \times L_m)}{\mathbf{Z}_m}
= \begin{bmatrix} \mathbf{z}_{1m}' \\ \vdots \\ \mathbf{z}_{nm}' \end{bmatrix}, \quad
\underset{(Mn \times \sum_{m=1}^{M} L_m)}{\mathbf{Z}}
= \begin{bmatrix}
\mathbf{Z}_1 & & \\
& \ddots & \\
& & \mathbf{Z}_M
\end{bmatrix},
\]
\[
\underset{(\sum_{m=1}^{M} L_m \times 1)}{\boldsymbol{\delta}}
= \begin{bmatrix} \boldsymbol{\delta}_1 \\ \vdots \\ \boldsymbol{\delta}_M \end{bmatrix}, \quad
\underset{(Mn \times 1)}{\mathbf{y}}
= \begin{bmatrix} \mathbf{y}_1 \\ \vdots \\ \mathbf{y}_M \end{bmatrix}, \quad
\underset{(n \times 1)}{\mathbf{y}_m}
= \begin{bmatrix} y_{1m} \\ \vdots \\ y_{nm} \end{bmatrix},
\]
\[
\underset{(Mn \times 1)}{\boldsymbol{\varepsilon}}
= \begin{bmatrix} \boldsymbol{\varepsilon}_1 \\ \vdots \\ \boldsymbol{\varepsilon}_M \end{bmatrix}, \quad
\underset{(n \times 1)}{\boldsymbol{\varepsilon}_m}
= \begin{bmatrix} \varepsilon_{1m} \\ \vdots \\ \varepsilon_{nm} \end{bmatrix}.
\]

Show the following:
\begin{itemize}
\item Assumption~4.1 becomes $\mathbf{y} = \mathbf{Z}\boldsymbol{\delta} + \boldsymbol{\varepsilon}$,
\item $\widehat{\mathbf{S}}$ in (4.5.10) becomes
$\widehat{\mathbf{S}} = \widehat{\boldsymbol{\Sigma}} \otimes \frac{1}{n} \mathbf{X}'\mathbf{X}$,
\item $\hat{\boldsymbol{\delta}}_{\mathrm{3SLS}}
= [\mathbf{Z}'(\widehat{\boldsymbol{\Sigma}}^{-1} \otimes \mathbf{P})\mathbf{Z}]^{-1}
\mathbf{Z}'(\widehat{\boldsymbol{\Sigma}}^{-1} \otimes \mathbf{P})\mathbf{y}$,
where $\mathbf{P} \equiv \mathbf{X}(\mathbf{X}'\mathbf{X})^{-1}\mathbf{X}'$,
\item $\widehat{\avar}(\hat{\boldsymbol{\delta}}_{\mathrm{3SLS}})
= n \cdot [\mathbf{Z}'(\widehat{\boldsymbol{\Sigma}}^{-1} \otimes \mathbf{P})\mathbf{Z}]^{-1}$,
\item $\widehat{\mathbf{A}}_{mh}$ in (4.5.13) becomes
$\widehat{\mathbf{A}}_{mh} = \frac{1}{n} \mathbf{Z}_m' \mathbf{P} \mathbf{Z}_h$,
\item $\hat{\mathbf{c}}_{mh}$ in (4.5.14) becomes
$\hat{\mathbf{c}}_{mh} = \frac{1}{n} \mathbf{Z}_m' \mathbf{P} \mathbf{y}_h$,
\item $J(\hat{\boldsymbol{\delta}}_{\mathrm{3SLS}},\, \widehat{\mathbf{S}}^{-1})
= (\mathbf{y} - \mathbf{Z}\hat{\boldsymbol{\delta}}_{\mathrm{3SLS}})'
(\widehat{\boldsymbol{\Sigma}}^{-1} \otimes \mathbf{P})
(\mathbf{y} - \mathbf{Z}\hat{\boldsymbol{\delta}}_{\mathrm{3SLS}})$.
\end{itemize}

\textbf{Hint:}
$\frac{1}{n} \sum_{i=1}^{n} \mathbf{x}_i \mathbf{x}_i' = \frac{1}{n} \mathbf{X}'\mathbf{X}$,
\quad $\mathbf{s}_{\mathbf{xy}} = \frac{1}{n}(\mathbf{I}_M \otimes \mathbf{X})'\mathbf{y}$,
\quad $\mathbf{S}_{\mathbf{xz}} = \frac{1}{n}(\mathbf{I}_M \otimes \mathbf{X})'\mathbf{Z}$.
Use formulas (A.11) and (A.13)--(A.15) of Appendix~A.

\exercise[2]
(Data matrix representation of SUR, RE, and Pooled OLS)
\begin{enumerate}[label=(\alph*)]
\item Prove that (4.5.13) reduces to (4.5.13$'$) on page~280 and (4.5.14) to
(4.5.14$'$), under the SUR assumption (4.5.18).
\textbf{Hint:} If $\mathbf{P} \equiv \mathbf{X}(\mathbf{X}'\mathbf{X})^{-1}\mathbf{X}'$,
then $\widehat{\mathbf{A}}_{mh} = \frac{1}{n} \mathbf{Z}_m' \mathbf{P} \mathbf{Z}_h$,
$\hat{\mathbf{c}}_{mh} = \frac{1}{n} \mathbf{Z}_m' \mathbf{P} \mathbf{y}_h$.
$\mathbf{P}\mathbf{Z}_m = \mathbf{Z}_m$ if the columns of $\mathbf{Z}_m$
are taken from the columns of $\mathbf{X}$.

\item Show that, for SUR,
\begin{align*}
\text{the estimator} &= [\mathbf{Z}'(\widehat{\boldsymbol{\Sigma}}^{-1} \otimes
\mathbf{I}_n)\mathbf{Z}]^{-1}
\mathbf{Z}'(\widehat{\boldsymbol{\Sigma}}^{-1} \otimes \mathbf{I}_n)\mathbf{y}, \\
\text{its } \widehat{\avar} &= n \cdot
[\mathbf{Z}'(\widehat{\boldsymbol{\Sigma}}^{-1} \otimes \mathbf{I}_n)\mathbf{Z}]^{-1}.
\end{align*}

\item Suppose, just for this part of the question, that $\mathbf{z}_{im} = \mathbf{x}_i$
for all $m$. So $\mathbf{Z} = \mathbf{I}_M \otimes \mathbf{X}$.
Show that SUR reduces to the multivariate regression estimator
\[
[\mathbf{I}_M \otimes (\mathbf{X}'\mathbf{X})^{-1}\mathbf{X}']\mathbf{y}.
\]
Verify that the $m$-th $K$-dimensional block of this $MK$-dimensional vector is
$(\mathbf{X}'\mathbf{X})^{-1}\mathbf{X}'\mathbf{y}_m$.

Now redefine $\mathbf{Z}$ as
\[
\underset{(Mn \times L)}{\mathbf{Z}}
\equiv \begin{bmatrix} \mathbf{Z}_1 \\ \vdots \\ \mathbf{Z}_M \end{bmatrix}. \qquad (*)
\]
So Assumption~4.1$'$ (linear equations with common coefficients) can be written
as $\mathbf{y} = \mathbf{Z}\boldsymbol{\delta} + \boldsymbol{\varepsilon}$,
where $\mathbf{y}$ and $\boldsymbol{\varepsilon}$ are as in Analytical Exercise~1
and $\boldsymbol{\delta}$ is the $L$-dimensional vector of common coefficients.

\item Show that, for the random-effects estimator, the same formulas you derived
in~(b) for SUR hold with $\mathbf{Z}$ as just redefined as in $(*)$.
\textbf{Hint:} The estimator is given in (4.6.8) and its $\widehat{\avar}$ is in
(4.6.10). Use formula (A.8) of the Appendix.

\item Show that
\begin{align*}
\hat{\boldsymbol{\delta}}_{\text{pooled OLS}} \text{ in (4.6.11)}
&= (\mathbf{Z}'\mathbf{Z})^{-1}\mathbf{Z}'\mathbf{y}, \\
\widehat{\avar}(\hat{\boldsymbol{\delta}}_{\text{pooled OLS}}) \text{ in (4.6.14)}
&= n \cdot (\mathbf{Z}'\mathbf{Z})^{-1}
[\mathbf{Z}'(\widehat{\boldsymbol{\Sigma}} \otimes \mathbf{I}_n)\mathbf{Z}]
(\mathbf{Z}'\mathbf{Z})^{-1}.
\end{align*}

\item In the $Mn \times L$ matrix $\mathbf{Z}$ defined in $(*)$, its rows are ordered
first by the equation $(m)$ and then by observation $(i)$. Now order the rows
first by observation $(i)$, which amounts to redefining $\mathbf{Z}$ as
\[
\underset{(Mn \times L)}{\mathbf{Z}}
= \begin{bmatrix} \mathbf{Z}_1 \\ \vdots \\ \mathbf{Z}_n \end{bmatrix},
\]
where $\mathbf{Z}_i$ $(i = 1, 2, \ldots, n)$ is the $M \times L$ matrix defined
in~(4.6.15). Reorder rows of $\mathbf{y}$ and $\boldsymbol{\varepsilon}$ similarly:
\[
\underset{(Mn \times 1)}{\mathbf{y}}
= \begin{bmatrix} \mathbf{y}_1 \\ \vdots \\ \mathbf{y}_n \end{bmatrix}, \quad
\underset{(Mn \times 1)}{\boldsymbol{\varepsilon}}
= \begin{bmatrix} \boldsymbol{\varepsilon}_1 \\ \vdots \\ \boldsymbol{\varepsilon}_n \end{bmatrix}.
\]

How would the formulas for RE and pooled OLS you derived in (d) and (e) change?
\textbf{Hint:} Just replace ``$\widehat{\boldsymbol{\Sigma}} \otimes \mathbf{I}_n$''
by ``$\mathbf{I}_n \otimes \widehat{\boldsymbol{\Sigma}}$''.
Set $M = 2$ if you find the translation too hard.
\end{enumerate}

\exercise[3]
(GLS interpretation of SUR and RE)\quad
The SUR model consists of Assumptions~4.1--4.5, 4.7, and (4.5.18)
(that $\mathbf{x}_i = \text{union of}\; (\mathbf{z}_{i1}, \ldots, \mathbf{z}_{iM})$).
Specialize the model by strengthening Assumption~4.3 (orthogonality conditions) as
\[
\E(\varepsilon_{im} \mid \mathbf{z}_{i1}, \ldots, \mathbf{z}_{iM}) = 0
\quad (m = 1, \ldots, M),
\]
and Assumption~4.2 (ergodic stationarity) as stating that
\[
\{y_{i1}, \ldots, y_{iM}, \mathbf{z}_{i1}, \ldots, \mathbf{z}_{iM}\}
\]
is i.i.d.

\begin{enumerate}[label=(\alph*)]
\item (Easy)\quad Explain why these are stronger assumptions.

\item Let $\boldsymbol{\varepsilon}$ $(Mn \times 1)$ and $\mathbf{Z}$
$(Mn \times \sum_m L_m)$ be as in Analytical Exercise~1 above.
Show that the specialized model implies
\[
\E(\boldsymbol{\varepsilon} \mid \mathbf{Z}) = \mathbf{0}
\quad \text{and} \quad
\E(\boldsymbol{\varepsilon}\boldsymbol{\varepsilon}' \mid \mathbf{Z})
= \boldsymbol{\Sigma} \otimes \mathbf{I}_n.
\]
\textbf{Hint:} For the first, what needs to be shown is
$\E(\varepsilon_{im} \mid \mathbf{Z}) = 0$. For the latter, the $(m, h)$ block of
$\E(\boldsymbol{\varepsilon}\boldsymbol{\varepsilon}' \mid \mathbf{Z})$ is
$\E(\boldsymbol{\varepsilon}_m \boldsymbol{\varepsilon}_h' \mid \mathbf{Z})$,
an $n \times n$ matrix. What needs to be shown is that this matrix is
$\sigma_{mh} \mathbf{I}_n$.

\item (Finite-sample properties of SUR)\quad
Suppose $\boldsymbol{\Sigma}$ is known so that the consistent estimates
$\hat{\sigma}_{mh}$ in the SUR formulas in Proposition~4.6 can be equated with
the true value $\sigma_{mh}$.
\begin{enumerate}[label=(\roman*)]
\item Verify that the model satisfies the GLS assumptions of Section~1.6
(which are Assumptions~1.1--1.3 and (1.6.1)) with a sample size of $Mn$.
\item Verify that the SUR estimator is the GLS estimator.
\textbf{Hint:} The $\mathbf{X}$ in Section~1.6 is $\mathbf{Z}$ here.
Show that $\sigma^2 \mathbf{V}$ there is $\boldsymbol{\Sigma} \otimes \mathbf{I}_n$.
Use your answer to 2(b).
\item Verify SUR is unbiased in finite samples in that
$\E(\hat{\boldsymbol{\delta}}_{\mathrm{SUR}} - \boldsymbol{\delta} \mid \mathbf{Z})
= \mathbf{0}$.
\item The expression for $\Var(\hat{\boldsymbol{\delta}}_{\mathrm{SUR}} \mid \mathbf{Z})$
is given in Proposition~1.7 of Section~1.6. Show that
\[
\avar(\hat{\boldsymbol{\delta}}_{\mathrm{SUR}})
= \plim_{n \to \infty} n \cdot
\Var(\hat{\boldsymbol{\delta}}_{\mathrm{SUR}} \mid \mathbf{Z}).
\]
\textbf{Hint:} With $\E(\boldsymbol{\varepsilon}\boldsymbol{\varepsilon}' \mid \mathbf{Z})
= \boldsymbol{\Sigma} \otimes \mathbf{I}_n$, the matrix
$\Var(\hat{\boldsymbol{\delta}}_{\mathrm{SUR}} \mid \mathbf{Z})$ is the inverse of a
matrix whose $(m,h)$ block is
$\sigma^{mh} \mathbf{Z}_m' \mathbf{Z}_h
= \sum_i \mathbf{z}_{im} \mathbf{z}_{ih}'$.
\end{enumerate}

\item (Finite-sample properties of RE)\quad
Now impose on the model the condition that the coefficient vector is the same
across equations so that $\mathbf{Z}$ is as in $(*)$ of Analytical Exercise~2.
As in (c), suppose $\boldsymbol{\Sigma}$ is known so that the consistent estimates
$\hat{\sigma}_{mh}$ in the RE formulas in Proposition~4.7 can be equated with
the true value $\sigma_{mh}$.
\begin{enumerate}[label=(\roman*)]
\item Verify that this model with common coefficients satisfies the GLS
assumptions of Section~1.6. In particular, Assumption~1.1 can be written as
$\mathbf{y} = \mathbf{Z}\boldsymbol{\delta} + \boldsymbol{\varepsilon}$,
where $\mathbf{y}$ and $\boldsymbol{\varepsilon}$ are as in Analytical Exercise~2.
\item Verify that the random-effects estimator is the GLS estimator.
\item Verify that the random-effects estimator is unbiased in finite samples.
\item Show that
\[
\avar(\hat{\boldsymbol{\delta}}_{\mathrm{RE}})
= \plim_{n \to \infty} n \cdot
\Var(\hat{\boldsymbol{\delta}}_{\mathrm{RE}} \mid \mathbf{Z}).
\]
\textbf{Hint:}
$\Var(\hat{\boldsymbol{\delta}}_{\mathrm{RE}} \mid \mathbf{Z})$ is the inverse of
$\sum_m \sum_h \sigma^{mh} \mathbf{Z}_m' \mathbf{Z}_h$.
\end{enumerate}
\end{enumerate}

\exercise[4]
(Standard errors from equation-by-equation vs.\ multiple-equation GMM)\quad
We have shown in Section~4.4 that the Avar of the efficient equation-by-equation
GMM estimator is no less (in matrix sense) than the Avar of the efficient
multiple-equation GMM estimator. In this exercise we show that the same
holds for $\widehat{\avar}$. Let $\widehat{\mathbf{S}}$ and
$\mathbf{S}_{\mathbf{xz}}$ be as in (4.3.2) and (4.2.2), respectively.

\begin{enumerate}[label=(\alph*)]
\item (Trivial)\quad Verify that, under appropriate assumptions,
$(\mathbf{S}_{\mathbf{xz}}' \widehat{\mathbf{S}}^{-1} \mathbf{S}_{\mathbf{xz}})^{-1}$
is consistent for the asymptotic variance of the efficient multiple-equation GMM
estimator. What are those appropriate assumptions?

\item Let $\tilde{\boldsymbol{\delta}}_m$ be the efficient single-equation GMM estimator
of $\boldsymbol{\delta}_m$ and $\tilde{\boldsymbol{\delta}}$ be the stacked vector
formed from $(\tilde{\boldsymbol{\delta}}_1, \ldots, \tilde{\boldsymbol{\delta}}_M)$.
$\tilde{\boldsymbol{\delta}}$ is the efficient equation-by-equation GMM estimator of
$\boldsymbol{\delta}$. Show that $\avar(\tilde{\boldsymbol{\delta}}_m)$ is consistently
estimated by
\[
\Bigl(\frac{1}{n} \sum_{i=1}^{n} \mathbf{z}_{im} \mathbf{x}_{im}'
\hat{\mathbf{W}}_{mm}
\frac{1}{n} \sum_{i=1}^{n} \mathbf{x}_{im} \mathbf{z}_{im}'\Bigr)^{-1},
\qquad (*)
\]
where
\[
\hat{\mathbf{W}}_{mm} = \Bigl(\frac{1}{n} \sum_{i=1}^{n}
\hat{\varepsilon}_{im}^2\, \mathbf{x}_{im} \mathbf{x}_{im}'\Bigr)^{-1}.
\]

\item Verify that $(*)$ is the $(m, m)$ block of
\[
(\mathbf{S}_{\mathbf{xz}}' \hat{\mathbf{W}} \widehat{\mathbf{S}}
\hat{\mathbf{W}} \mathbf{S}_{\mathbf{xz}})
(\mathbf{S}_{\mathbf{xz}}' \hat{\mathbf{W}} \mathbf{S}_{\mathbf{xz}})^{-1},
\qquad (**)
\]
where $\hat{\mathbf{W}}$ is the block diagonal matrix whose $m$-th diagonal block
is $\hat{\mathbf{W}}_{mm}$.

\item Proposition~3.5 is an algebraic result applicable not just to population
moments, so it implies that
\[
(**) \geq (\mathbf{S}_{\mathbf{xz}}' \widehat{\mathbf{S}}^{-1}
\mathbf{S}_{\mathbf{xz}})^{-1}.
\]
Taking this for granted, show the following: If the same residuals are used
to calculate $\hat{\mathbf{W}}$ for the equation-by-equation GMM and
$\widehat{\mathbf{S}}$ for the multiple-equation GMM, then, for each coefficient,
the standard error from the efficient multiple-equation GMM is no larger than
that from the efficient equation-by-equation GMM.
\textbf{Hint:} The standard errors are the square roots of the diagonal elements of
the estimated asymptotic variance divided by the sample size.

\item Does the result you proved in (d) hold true if the model is misspecified in
any way? \textbf{Hint:} What you have proved is an algebraic result.

\item Go through steps (a)--(d) for the FIVE and equation-by-equation 2SLS.
\textbf{Hint:} It's just a matter of redefining $\widehat{\mathbf{S}}$ and the $m$-th
diagonal block, $\hat{\mathbf{W}}_{mm}$, of $\hat{\mathbf{W}}$.
The $\hat{\mathbf{W}}_{mm}$ for 2SLS is the inverse of
\[
\frac{1}{n} \sum_{i=1}^{n} \hat{\varepsilon}_{im}^2 \cdot
\frac{1}{n} \sum_{i=1}^{n} \mathbf{x}_{im} \mathbf{x}_{im}'
\]
and the $\widehat{\mathbf{S}}$ for FIVE is (4.5.3).
\end{enumerate}

\exercise[5]
(Identification and cross-equation restrictions)\quad
Consider the wage equation for 1969 and 1980:
\begin{align*}
LW69_i &= \beta_0 + \beta_1 \cdot S69_i + \beta_2 \cdot IQ_i + \varepsilon_{i1}, \\
LW80_i &= \beta_0 + \beta_1 \cdot S80_i + \beta_2 \cdot IQ_i + \varepsilon_{i2}.
\end{align*}
The orthogonality conditions are $\E(\varepsilon_{i1}) = 0$,
$\E(\varepsilon_{i2}) = 0$,
$\E(MED_i\, \varepsilon_{i1}) = 0$.
Assume we have a random sample on $(LW69, LW80, S69, S80, IQ, MED)$.

\begin{enumerate}[label=(\alph*)]
\item Is the $LW69$ equation identified in isolation?
\textbf{Hint:} Check the order condition for identification for the $LW69$ equation.

\item Write the four orthogonality conditions in terms of
$(\beta_0, \beta_1, \beta_2)$ as a system of equations linear in $\beta$'s.
How many equations are there? State the rank condition for identification in terms
of relevant cross moments. (This example shows that even if each equation of the
system is not identified in isolation, the system can be identified thanks to the
cross-equation restriction that $\beta$'s are the same across equations.)

\item Show that the system cannot be identified if $IQ$ and $MED$ are uncorrelated.
\end{enumerate}

\exercise[6]
(Optional)\quad Prove Proposition~4.1. \textbf{Hint:} Generalize the proof of
Proposition~3.2.

\exercise[7]
(Numerical equivalence of 2SLS and 3SLS)\quad
In the system of equations
$y_{im} = \mathbf{z}_{im}' \boldsymbol{\delta}_m + \varepsilon_{im}$
$(i = 1, \ldots, n;\; m = 1, \ldots, M)$, suppose all equations share the
same set of instruments (so $\mathbf{x}_{im} = \mathbf{x}_i$). (In what follows,
assume all the conditions for the 3SLS estimator are satisfied.)

\begin{enumerate}[label=(\alph*)]
\item When each equation has the same RHS variables (so $\mathbf{z}_{im} = \mathbf{z}_i$),
the 2SLS and 3SLS estimators are numerically the same. Prove this.
\textbf{Hint:} Let $\mathbf{B} \equiv \frac{1}{n} \sum_i \mathbf{x}_i \mathbf{z}_i'$.
The 3SLS estimator is (4.2.3) with $\mathbf{S}_{\mathbf{xz}} = \mathbf{I}_M \otimes \mathbf{B}$,
$\hat{\mathbf{W}} = \widehat{\boldsymbol{\Sigma}}^{-1} \otimes \mathbf{S}_{\mathbf{xx}}^{-1}$
where $\mathbf{S}_{\mathbf{xx}} \equiv \frac{1}{n} \sum_i \mathbf{x}_i \mathbf{x}_i'$.

\item If the errors are not conditionally homoskedastic, does (a) remain true for
the single-equation and multiple-equation GMM estimators?
\textbf{Hint:} The answer is no. Why? Can you write $\widehat{\mathbf{S}}^{-1}$ as a
Kronecker product without conditional homoskedasticity?
\end{enumerate}

\exercise[8]
(SUR is more than assuming predetermined regressors in each equation)\quad
Consider estimating a system of $M$ equations:
$y_{im} = \mathbf{z}_{im}' \boldsymbol{\delta}_m + \varepsilon_{im}$
$(i = 1, \ldots, n;\; m = 1, \ldots, M)$.
The orthogonality conditions are
$\E(\mathbf{z}_{im} \cdot \varepsilon_{im}) = \mathbf{0}$
$(m = 1, 2, \ldots, M)$.

\begin{enumerate}[label=(\alph*)]
\item Does the efficient multiple-equation GMM estimator
$\hat{\boldsymbol{\delta}}(\widehat{\mathbf{S}}^{-1})$ reduce to the
equation-by-equation OLS? Does your conclusion change if the errors are
conditionally homoskedastic?
\textbf{Hint:} What's the size of $\mathbf{S}_{\mathbf{xz}}$? $\widehat{\mathbf{S}}$?

\item Explain why the efficient GMM you derived in (a) under conditional
homoskedasticity differs from SUR.
\textbf{Hint:} The difference is in the orthogonality conditions.
\end{enumerate}

\exercise[9]
(Optional, peril of joint estimation)\quad
Taken from Greene (1997, p.~706). Consider the two-equation SUR model:
\begin{align*}
y_{i1} &= \beta_1 x_{i1} + \varepsilon_{i1}, \\
y_{i2} &= \beta_2 x_{i2} + \beta_3 x_{i3} + \varepsilon_{i2}.
\end{align*}
For simplicity, assume that $\E(\varepsilon_{i1}^2)$,
$\E(\varepsilon_{i1}\, \varepsilon_{i2})$, and $\E(\varepsilon_{i2}^2)$ are known.
Now suppose you apply SUR but erroneously omit $x_{i3}$ from the second equation.
What effect does this have on the consistency of the estimator of $\beta_1$?
\textbf{Hint:} Let $\mathbf{x}_1$, $\mathbf{x}_2$, $\mathbf{x}_3$,
$\mathbf{y}_1$, $\mathbf{y}_2$, $\boldsymbol{\varepsilon}_1$, and
$\boldsymbol{\varepsilon}_2$ be $n$-dimensional vectors of respective variables and
\[
\mathbf{d} = \begin{bmatrix} \mathbf{0} \\ \mathbf{x}_3 \end{bmatrix}, \quad
\boldsymbol{\beta} = \begin{bmatrix} \beta_1 \\ \beta_2 \end{bmatrix}, \quad
\boldsymbol{\Sigma} = \begin{bmatrix}
\E(\varepsilon_{i1}^2) & \E(\varepsilon_{i1}\, \varepsilon_{i2}) \\
\E(\varepsilon_{i2}\, \varepsilon_{i1}) & \E(\varepsilon_{i2}^2)
\end{bmatrix}.
\]
(So, for example, $\bar{\mathbf{X}}$ is $2n \times 2$ and $\boldsymbol{\beta}$ is
$2 \times 1$.) Show that:
\begin{enumerate}[label=(\alph*)]
\item The two-equation system can be written as
$\mathbf{y} = \bar{\mathbf{X}}\boldsymbol{\beta} + \beta_3 \cdot \mathbf{d}
+ \boldsymbol{\varepsilon}$.

\item The SUR estimator of $\boldsymbol{\beta}$ when $x_{i3}$ is inadvertently left
out is
\[
\hat{\boldsymbol{\beta}}_{\mathrm{SUR}}
= [\bar{\mathbf{X}}'(\boldsymbol{\Sigma}^{-1} \otimes \mathbf{I}_n)
\bar{\mathbf{X}}]^{-1}
\bar{\mathbf{X}}'(\boldsymbol{\Sigma}^{-1} \otimes \mathbf{I}_n)\mathbf{y}.
\]

\item Use formulas (A.6) and (A.7) of the Appendix to show
$\hat{\boldsymbol{\beta}}_{\mathrm{SUR}} = \boldsymbol{\beta} + \beta_3 \cdot \mathbf{a}
+ \mathbf{b}$,
where $\mathbf{a}$ and $\mathbf{b}$ are given explicitly.

\item Argue that $\plim \mathbf{b} = \mathbf{0}$. Is $\plim \mathbf{a} = \mathbf{0}$?
\end{enumerate}

\exercise[10]
(Optional, adding just-identified equations does not increase efficiency)\quad
Consider the 3SLS two-equation model where the second equation is just identified
(so $L_2 = K$). Let $\hat{\boldsymbol{\delta}}_{1,\mathrm{3SLS}}$ be the 3SLS
estimator of the coefficients of the first equation and
$\hat{\boldsymbol{\delta}}_{1,\mathrm{2SLS}}$ be the 2SLS estimator of the same
coefficients.

\begin{enumerate}[label=(\alph*)]
\item Write down $\avar(\hat{\boldsymbol{\delta}}_{1,\mathrm{2SLS}})$ in terms of
$\sigma_{11}$ and $\mathbf{A}_{11}$, which is defined in (4.5.16).

\item Show that
$\avar(\hat{\boldsymbol{\delta}}_{1,\mathrm{3SLS}})
= \avar(\hat{\boldsymbol{\delta}}_{1,\mathrm{2SLS}})$.
\textbf{Hint:} Use the partitioned inverse formula (see formula (A.10) of the
Appendix) to write the upper-left ($L_1 \times L_1$) block of
$\avar(\hat{\boldsymbol{\delta}}_{1,\mathrm{3SLS}})$ as $\mathbf{G}^{-1}$ where
\[
\mathbf{G} = \sigma^{11} \mathbf{A}_{11}
- \Bigl(\frac{\sigma^{12} \cdot \sigma^{21}}{\sigma^{22}}\Bigr)
\mathbf{A}_{12} \mathbf{A}_{22}^{-1} \mathbf{A}_{21}.
\]
But, by (4.5.16),
\begin{align*}
\mathbf{A}_{12} \mathbf{A}_{22}^{-1} \mathbf{A}_{21}
&= \E(\mathbf{z}_{i1} \mathbf{x}_i')
[\E(\mathbf{x}_i \mathbf{x}_i')]^{-1}
\E(\mathbf{x}_i \mathbf{z}_{i2}') \\
&\quad \times [\E(\mathbf{z}_{i2} \mathbf{x}_i')
[\E(\mathbf{x}_i \mathbf{x}_i')]^{-1}
\E(\mathbf{x}_i \mathbf{z}_{i2}')]^{-1} \\
&\quad \times \E(\mathbf{z}_{i2} \mathbf{x}_i')
[\E(\mathbf{x}_i \mathbf{x}_i')]^{-1}
\E(\mathbf{x}_i \mathbf{z}_{i1}') \\
&= \E(\mathbf{z}_{i1} \mathbf{x}_i')
[\E(\mathbf{x}_i \mathbf{x}_i')]^{-1}
\E(\mathbf{x}_i \mathbf{z}_{i1}')
\quad (\text{since } \#\mathbf{z}_{i2} (= L_2) = K (= \#\mathbf{x}_i)) \\
&= \mathbf{A}_{11}.
\end{align*}
So
\[
\mathbf{G} = \left[\sigma^{11}
- \Bigl(\frac{\sigma^{12} \cdot \sigma^{21}}{\sigma^{22}}\Bigr)\right]
\mathbf{A}_{11}
= \frac{\mathbf{A}_{11}}{\sigma_{11}}
\quad \bigl(\text{since } \sigma^{11}
- \sigma^{12}\sigma^{21}/\sigma^{22} = 1/\sigma_{11}\bigr).
\]
\end{enumerate}

\exercise[11]
(Linear combinations of orthogonality conditions)\quad
For the multiple-equation model of Proposition~4.7, derive the efficient GMM
estimator that exploits the following orthogonality conditions:
\begin{equation}
\E(\mathbf{z}_{i1} \cdot \varepsilon_{i1} + \cdots
+ \mathbf{z}_{iM} \cdot \varepsilon_{iM}) = \mathbf{0}.
\tag{1}
\end{equation}
\textbf{Hint:} The sample analogue of the left-hand side of the orthogonality conditions is
\[
\mathbf{g}_n(\tilde{\boldsymbol{\delta}})
\equiv \frac{1}{n} \sum_{i=1}^{n} \mathbf{z}_{i1} \cdot y_{i1}
+ \cdots + \frac{1}{n} \sum_{i=1}^{n} \mathbf{z}_{iM} \cdot y_{iM}
- \Bigl(\frac{1}{n} \sum \mathbf{z}_{i1} \mathbf{z}_{i1}'
+ \cdots + \frac{1}{n} \sum \mathbf{z}_{iM} \mathbf{z}_{iM}'\Bigr)
\tilde{\boldsymbol{\delta}}.
\]
There are as many orthogonality conditions as there are coefficients.
The efficient GMM estimator is the IV estimator that solves
$\mathbf{g}_n(\tilde{\boldsymbol{\delta}}) = \mathbf{0}$.

\bigskip
\subsection*{Empirical Exercises}

In data file \texttt{GREENE.ASC}, data are provided on the following variables.

\medskip
\begin{tabular}{@{}ll@{}}
Column 1: & firm ID in Nerlove's 1955 sample \\
Column 2: & costs in 1970 in millions of dollars \\
Column 3: & output in millions of kilowatt hours \\
Column 4: & price of labor \\
Column 5: & price of capital \\
Column 6: & price of fuel \\
Column 7: & labor's cost share \\
Column 8: & capital's cost share
\end{tabular}
\medskip

The sample has 99 observations. Fuel's cost share can be calculated as one minus
the sum of labor and capital cost shares. These data were used in Christensen and
Greene (1976).

In the text, the third equation (for fuel share) of the system was dropped in
the random-effects estimation. To verify the numerical invariance, drop the second
equation (for capital share) and eliminate $(\gamma_{12}, \gamma_{22}, \gamma_{32})$
when imposing homogeneity. This produces the two-equation system:
\[
\begin{cases}
s_1 = \alpha_1 + \gamma_{11} \log(p_1/p_2) + \gamma_{13} \log(p_3/p_2)
+ \gamma_{1Q} \log(Q) + \varepsilon_1
\quad \text{(labor share),} \\
s_3 = \alpha_3 + \gamma_{31} \log(p_1/p_2) + \gamma_{33} \log(p_3/p_2)
+ \gamma_{3Q} \log(Q) + \varepsilon_3
\quad \text{(fuel share).}
\end{cases}
\]
Call this the \textit{unconstrained system}. The \textit{constrained system} imposes
symmetry $\gamma_{13} = \gamma_{31}$ on the unconstrained system.

\begin{enumerate}[label=(\alph*)]
\item Verify that the simple statistics of the sample shown in the text can be
duplicated.

\item (Numerical invariance)\quad
Apply the random-effects estimation (i.e., the multivariate regression subject to
the cross-equation restriction) to the constrained system. Can you duplicate the
coefficient estimates in the text? You should first apply equation-by-equation OLS
to the unconstrained system to calculate
$\widehat{\boldsymbol{\Sigma}}^*$. (Note on computation: Programming joint
estimation is much easier with canned programs such as TSP and RATS.
However, those canned programs will not print out Sargan's statistic for the SUR.)

\textit{Gauss Tip:} The random-effects estimator could be calculated by operating
on data matrices, but we recommend the use of do loops.
Use (4.6.8$'$) on page~293.

\textit{RATS Tip:} Use the \texttt{NLSYS} proc, which is designed for handling nonlinear
equations but which can be used for linear equations. It can also be used for
calculating $\widehat{\boldsymbol{\Sigma}}^*$.

\textit{TSP Tip:} See codes in the text for accomplishing the same.

This gives the random-effects estimate of the seven free parameters,
\[
\alpha_1, \alpha_3, \gamma_{11}, \gamma_{13}, \gamma_{33},
\gamma_{1Q}, \gamma_{3Q}.
\]
As explained in the text, use the adding-up, homogeneity, and symmetry
restrictions to calculate the point estimates and their standard errors of the
remaining eight parameters.

\item (Optional, for Gauss users)\quad Calculate Sargan's statistic.
\textbf{Hint:} Use the formula in Proposition~4.7.
The $\mathbf{x}_i$ vector should be: a constant, $\log(p_1/p_2)$,
$\log(p_3/p_2)$, $\log(Q)$.
For $\widehat{\boldsymbol{\Sigma}}^*$, use the one obtained from the unconstrained
system.
$\mathbf{s}_{\mathbf{xy}} = \frac{1}{n} \sum_{i=1}^{n} \mathbf{y}_i \otimes \mathbf{x}_i$
and $\mathbf{S}_{\mathbf{xz}} = \frac{1}{n} \sum_{i=1}^{n} \mathbf{Z}_i \otimes \mathbf{x}_i$.
The statistic should be 0.63313 with a $p$-value of 0.42621.

\item (Testing symmetry)\quad
Given that Sargan's statistic is 0.63313, perform the following hypothesis testing.
Take the maintained hypothesis to be the unconstrained system.
The null hypothesis is the symmetry restriction that $\gamma_{13} = \gamma_{31}$.
\textbf{Hint:} By Proposition~3.8, the difference in the Sargan statistics with and
without symmetry is asymptotically chi-squared. What is Sargan's statistic for
the unconstrained system?

\item (Wald test of symmetry)\quad
Test the same null hypothesis by the Wald principle.
So you test the symmetry restriction in the unconstrained system. Verify that
the test statistic is numerically equal to Sargan's statistic for the constrained
system you estimated in (b).

\item Using the parameter estimates you obtained in (b), calculate the elasticity of
substitution between labor and capital averaged over the 99 firms. Can you
duplicate the result in the text?
\end{enumerate}

\end{problemset}

\begin{answers}

\subsection*{Analytical Exercises}

\answer{3b}
$\E(\varepsilon_{im} \mid \mathbf{Z})$
\begin{align*}
&= \E(\varepsilon_{im} \mid \mathbf{z}_{11}, \ldots, \mathbf{z}_{1M},
\mathbf{z}_{21}, \ldots, \mathbf{z}_{2M}, \ldots, \mathbf{z}_{n1}, \ldots, \mathbf{z}_{nM})
\quad \text{(by definition of $\mathbf{Z}$)} \\
&= \E(\varepsilon_{im} \mid \mathbf{z}_{i1}, \ldots, \mathbf{z}_{iM})
\quad \text{(by the i.i.d.\ assumption)} \\
&= 0 \quad \text{(by the strengthened orthogonality conditions).}
\end{align*}

The $(i, j)$ element of the $n \times n$ matrix
$\E(\boldsymbol{\varepsilon}_m \boldsymbol{\varepsilon}_h' \mid \mathbf{Z})$
is $\E(\varepsilon_{im}\, \varepsilon_{jh} \mid \mathbf{Z})$.
By the i.i.d.\ assumption, this is zero for $j \neq i$. For $j = i$,
\begin{align*}
\E(\varepsilon_{im}\, \varepsilon_{ih} \mid \mathbf{Z})
&= \E(\varepsilon_{im}\, \varepsilon_{ih} \mid
\mathbf{z}_{i1}, \ldots, \mathbf{z}_{1M}, \mathbf{z}_{21}, \ldots, \mathbf{z}_{nM}) \\
&= \E(\varepsilon_{im}\, \varepsilon_{ih} \mid \mathbf{z}_{i1}, \ldots, \mathbf{z}_{iM})
\quad \text{(by the i.i.d.\ assumption).}
\end{align*}
Since $\mathbf{x}_{im} = \mathbf{x}_i$ and $\mathbf{x}_i$ is the union of
$(\mathbf{z}_{i1}, \ldots, \mathbf{z}_{iM})$ in the SUR model, the conditional
homoskedasticity assumption, Assumption~4.7, states that
$\E(\varepsilon_{im}\, \varepsilon_{ih} \mid \mathbf{z}_{i1}, \ldots, \mathbf{z}_{iM})
= \sigma_{mh}$.

\answer{5b}
\[
\begin{bmatrix}
1 & \E(S69) & \E(IQ) \\
1 & \E(S80) & \E(IQ) \\
\E(MED) & \E(S69 \cdot MED) & \E(IQ \cdot MED) \\
\E(MED) & \E(S80 \cdot MED) & \E(IQ \cdot MED)
\end{bmatrix}
\begin{bmatrix} \beta_0 \\ \beta_1 \\ \beta_2 \end{bmatrix}
= \begin{bmatrix}
\E(LW69) \\ \E(LW80) \\ \E(LW69 \cdot MED) \\ \E(LW80 \cdot MED)
\end{bmatrix}.
\]
The condition for the system to be identified is that the $4 \times 3$ coefficient matrix
is of full column rank.

\answer{5c}
If $IQ$ and $MED$ are uncorrelated, then
$\E(IQ \cdot MED) = \E(IQ) \cdot \E(MED)$ and the third column of the coefficient
matrix is $\E(IQ)$ times the first column. So the matrix cannot be of full column rank.

\answer{6}
$\hat{\varepsilon}_{im} = y_{im} - \mathbf{z}_{im}' \hat{\boldsymbol{\delta}}_m
= \varepsilon_{im} - \mathbf{z}_{im}'(\hat{\boldsymbol{\delta}}_m - \boldsymbol{\delta}_m)$.
So
\[
\frac{1}{n} \sum_{i=1}^{n}
[\varepsilon_{im} - \mathbf{z}_{im}'(\hat{\boldsymbol{\delta}}_m - \boldsymbol{\delta}_m)]
[\varepsilon_{ih} - \mathbf{z}_{ih}'(\hat{\boldsymbol{\delta}}_h - \boldsymbol{\delta}_h)]
= (1) + (2) + (3) + (4),
\]
where
\begin{align*}
(1) &= \frac{1}{n} \sum \varepsilon_{im}\, \varepsilon_{ih}, \\
(2) &= -(\hat{\boldsymbol{\delta}}_m - \boldsymbol{\delta}_m)'
\frac{1}{n} \sum \mathbf{z}_{im} \cdot \varepsilon_{ih}, \\
(3) &= -(\hat{\boldsymbol{\delta}}_h - \boldsymbol{\delta}_h)'
\frac{1}{n} \sum \mathbf{z}_{ih} \cdot \varepsilon_{im}, \\
(4) &= (\hat{\boldsymbol{\delta}}_m - \boldsymbol{\delta}_m)'
\Bigl(\frac{1}{n} \sum \mathbf{z}_{im} \mathbf{z}_{ih}'\Bigr)
(\hat{\boldsymbol{\delta}}_h - \boldsymbol{\delta}_h).
\end{align*}

As usual, under Assumptions~4.1 and~4.2,
$(1) \pto \sigma_{mh}$ $(\equiv \E(\varepsilon_{im}\, \varepsilon_{ih}))$.
For $(4)$, by Assumption~4.2 and the assumption that
$\E(\mathbf{z}_{im} \mathbf{z}_{ih}')$ is finite,
$\frac{1}{n} \sum_i \mathbf{z}_{im} \mathbf{z}_{ih}'$ converges in probability to
a (finite) matrix. So $(4) \pto 0$.

Regarding $(2)$, by Cauchy-Schwartz,
\[
\E(|z_{imj} \cdot \varepsilon_{ih}|) \leq \sqrt{\E(z_{imj}^2) \cdot \E(\varepsilon_{ih}^2)},
\]
where $z_{imj}$ is the $j$-th element of $\mathbf{z}_{im}$.
So $\E(\mathbf{z}_{im} \cdot \varepsilon_{ih})$ is finite and
$(2) \pto 0$ because $\hat{\boldsymbol{\delta}}_m - \boldsymbol{\delta}_m \pto \mathbf{0}$.
Similarly, $(3) \pto 0$.
\end{answers}
